\section{System Design}\label{sec:design}

\begin{figure*}[t]
\centering
\includegraphics[width=0.85\textwidth]{pipeline}
\caption{Pipeline. The two highlighted stages are the only things that change between rows of a result table;
everything to their right is fixed.}
\label{fig:pipeline}
\end{figure*}

The research question (``change only the depth source---how much changes?'') leads to one design rule: every depth
source passes through the same code after the point where it is produced, and whatever can be swapped is swapped from
a YAML config, so one result table is one config file (Fig.~\ref{fig:pipeline}).

\textbf{DepthSource.} An interface \code{get\_depth(frame)} returning an $(H,W)$ depth map and a flag
\code{is\_metric}. Implementations read Faro depth (\code{gt}), ARKit LiDAR (\code{lidar}) or run a monocular network
(\code{mono}: DA-v2, MiDaS, fine-tuned checkpoints) on the upright RGB frame.

\textbf{ScaleAligner.} A relative model is correct only up to an affine map, and \emph{how} that map is obtained is a
system decision measured separately from the model (Table~\ref{tab:aligners}). Affine-invariant models are affine in
\emph{disparity}, so every aligner fits $1/gt \approx s\cdot(1/pred) + t$ by trimmed least squares and inverts back;
on the development scene this change alone moved mono$+$oracle from \SI{15.2}{cm} to \SI{5.4}{cm}. Metric sources
must use \code{identity}; the pipeline refuses to build otherwise.

\begin{table}[t]
\centering\small
\caption{Scale aligners. Every non-identity aligner fits $(s,t)$ in inverse depth with trimmed least squares.}
\label{tab:aligners}
\begin{tabularx}{\columnwidth}{@{}lX@{}}
\toprule
Aligner & Fitted on / meaning \\
\midrule
\code{identity} & nothing; for Faro, LiDAR and metric models \\
\code{oracle\_frame} & that frame's GT; the model's ceiling (not deployable) \\
\code{per\_scene} & GT of 10 frames spread over the scan, fitted once; upper bound of a one-off calibration \\
\code{sparse\_points} & 200 LiDAR pixels per frame at confidence 2, a proxy for VIO feature points; upper bound of a
VIO-based system \\
\bottomrule
\end{tabularx}
\end{table}

\textbf{Fixed stages.} Open3D \code{ScalableTSDFVolume} with \SI{4}{cm} voxels, truncation $3\times$ voxel and depth
range \SIrange{0.1}{5}{m}; every source is resampled to the LiDAR resolution ($256\times192$) for fairness, and only
pixels where Faro has depth are integrated, so no source is penalised for seeing beyond the reference. Marching cubes
and small-cluster removal follow. Poses are the dataset's ARKit VIO trajectory for every row.
